\documentclass[10pt,a4paper]{amsart}
\usepackage[margin=19mm]{geometry}
\usepackage{amsmath,amssymb,mathtools}
\usepackage[scr=rsfs,cal=euler]{mathalfa}
\usepackage{xfrac}
\usepackage[unicode,hidelinks,bookmarks=false]{hyperref}
\newcommand{\ie}{i.e.\ }
\newcommand{\cf}{cf.\ }
\newcommand{\resp}{respectively\ }
\newcommand{\Subsection}{\subsection}
\providecommand{\nobreakdash}{\nobreak\hbox{-}}
\setlength{\emergencystretch}{2em}
\multlinegap0pt
\usepackage{amssymb}
\newtheorem{assumption}{Assumption}
\newtheorem{theorem}{Theorem}
\newtheorem{lemma}{Lemma}
\newtheorem{definition}{Definition}
\newtheorem{corollary}{Corollary}
\def\Ga{\Gamma}
\def\La{\Lambda}
\def\Om{\Omega}
\def\be{\beta}
\def\df{\mathrm d}
	\DeclareMathOperator{\dist}{dist}
\def\esssup{\operatorname*{ess\ \! sup}}
\def\f{\frac}
\def\lg{\langle}
\def\mcA{\mathcal{A}}
\def\mcD{\mathcal{D}}
\def\mcO{\mathcal{O}}
\def\ol{\overline}
\def\pt{\partial}
\def\ra{\rightarrow}
\def\rg{\rangle}
\def\ts{\times}
\def\wh{\widehat}
\title{Non-homogeneous boundary value problems for second-order degenerate hyperbolic equations and their application}
\author{Donghui Yang, Jie Zhong}
\date{Source: arXiv:2602.07271v1 (CC BY 4.0)}
\begin{document}
\maketitle

\noindent\emph{Source and licence.} Non-homogeneous boundary value problems for second-order degenerate hyperbolic equations and their application. Version arXiv:2602.07271v1 (CC BY 4.0), distributed by arXiv under the Creative Commons Attribution 4.0 International licence, http://creativecommons.org/licenses/by/4.0/, as recorded in the arXiv licence metadata for this exact version and shown on the abstract page https://arxiv.org/abs/2602.07271v1. Full licence notice: \texttt{licenses/arxiv-2602.07271-CC-BY-4.0.txt}.\par\smallskip

\noindent\textit{Editorial scope.} Counted unit: the article's theorem 09.02.T1 (source0.tex lines 725--750). The article's complete original proof of it is delivered with it (source0.tex lines 903--973, one \texttt{proof} environment of 71 lines, which the article places away from the statement). No mathematics is written, repaired or re-derived: the statement and the proof are the article's own lines, in the article's own notation, and no retained line is substituted or re-flowed.  Retained uncounted as the source-local context the proof needs: lines 474--491; lines 537--546; lines 559--577; lines 687--699; lines 716--722; lines 764--773; lines 829--854; lines 861--870; lines 1364--1366; lines 1373--1380; lines 1428--1434.  Nothing was omitted from any retained span, so the package carries no omission record.  No retained line contains a citation, so the wrapper carries no bibliography and no bibliography entry.  No retained line contains a figure environment, an image-inclusion command or drawing code, so no image is delivered and the package declares no asset dependency.  Editorial formatting only, in addition: an AMS \texttt{amsart} preamble that restates the article's own declarations and macros used by the retained lines, namely the environments \texttt{assumption}, \texttt{theorem}, \texttt{lemma}, \texttt{definition}, \texttt{corollary} on separate counters, together with the standard packages those lines need.  The retained environments print in this document as Assumption 1, Theorem 1, Lemma 1, Definition 1, Corollary 1 in order of appearance, whereas the article prints the same items with the numbering of its own full document.  The complete native source of the article as distributed in the arXiv e-print for this exact version is bundled unchanged as \texttt{sources/194167/source0.tex}.
\begin{assumption}\label{06.27.A1}\label{ass:2.2}
In addition to $w\in A_2$, we assume:
\begin{enumerate}
  \item (\textit{Non-degeneracy near the boundary}) There exist constants $\Lambda>0$ and $\beta>0$ such that
	\begin{equation}\label{08.13.1}
		w\in C^1(\ol{\mcO(\Ga;2\be)}) \mbox{ and } w\geq \La \mbox{ on } \ol{\mcO(\Ga;2\beta)},
	\end{equation}
	where $\mcO(\Ga;\beta)=\left\{x\in\Om\colon \dist(x,\Ga)<\beta\right\}$ and $\dist(x,\Ga)=\inf_{z\in \Ga}|x-z|$.
  \item (\textit{Weighted Poincar\'e inequality}) There exists $C>0$ such that for all $y\in H^1_0(\Omega;w)$,
\[
\int_{\Omega} |y(x)|^2\, dx \le C_P \int_{\Omega} w(x)|\nabla y(x)|^2\, dx.
\]

  \item (\textit{Compact embedding}) $H^1_0(\Omega;w)$ is compactly embedded in $L^2(\Omega)$.

  \item (\textit{Initial trace compatibility}) For well-posedness of boundary input problems, we assume $u(0,\xi) = y_0(\xi) = 0$ for $\xi\in \Gamma$.
\end{enumerate}
\end{assumption}

We now consider the homogeneous Dirichlet problem for the degenerate wave equation,
\begin{equation}
\begin{cases}
\partial_{tt}y + Ay = f(t,x), & (x,t)\in Q=\Omega\times(0,T),\\
y(0,x) = y_0(x),\quad \partial_t y(0,x) = y_1(x), & x\in \Omega,\\
y(t,\xi) = 0, & (\xi,t)\in \Sigma=\Gamma\times(0,T),
\end{cases}
\label{09.17.1}
\end{equation}
with data $f\in L^1(0,T;L^2(\Omega))$, $y_0\in H^1_0(\Omega;w)$, and $y_1\in L^2(\Omega)$. The definition below is the notion of solution used in Theorem~\ref{thm:2.12}.

\begin{theorem}[Existence and regularity]\label{thm:2.12}
Under Assumption~\ref{ass:2.2}:
\begin{enumerate}
\item For $y_0\in H^1_0(\Omega;w)$, $y_1\in L^2(\Omega)$, and $f\in L^1(0,T;L^2(\Omega))$, there exists a unique weak solution $y$ of \eqref{09.17.1} such that
\[
\esssup_{0\le t\le T} \big(\|y(t)\|_{H^1_0(\Omega;w)} + \|\partial_t y(t)\|_{L^2(\Omega)}\big)
\le C\big(\|y_0\|_{H^1_0(\Omega;w)} + \|y_1\|_{L^2(\Omega)} + \|f\|_{L^1(0,T;L^2(\Omega))}\big).
\]
\item If $y_0\in D(A)$, $y_1\in H^1_0(\Omega;w)$, and $\partial_t f \in L^1(0,T;L^2(\Omega))$, then
\[
y \in L^\infty(0,T;D(A)) \cap W^{1,\infty}(0,T;H^1_0(\Omega;w)), \quad \partial_{tt}y \in L^\infty(0,T;L^2(\Omega)),
\]
and
\[
\esssup_{0\le t\le T}\Big(\|y(t)\|_{D(A)} + \|\partial_t y(t)\|_{H^1_0(\Omega;w)} + \|\partial_{tt}y(t)\|_{L^2(\Omega)}\Big)
\le C\big(\|y_0\|_{D(A)} + \|y_1\|_{H^1_0(\Omega;w)} + \|f\|_{L^1} + \|\partial_t f\|_{L^1}\big).
\]
\end{enumerate}
\end{theorem}

\begin{lemma}\label{09.11.L1}
	Under Assumption \ref{06.27.A1}, the degenerate elliptic equation \begin{equation}\label{09.11.1}
		\begin{cases}
			Av=f, &\mbox{in }\Om,\\
			v=\psi, &\mbox{on }\pt\Ga, 
		\end{cases}
	\end{equation}
	has a weak solution in $H^1(\Om;w)$, where $f\in H^{-1}(\Om;w)$ and $\psi\in H^1(\Ga)$. Moreover,  
	\begin{equation}\label{09.12.4}
		\|v\|_{H^1(\Om;w)}\leq C\left(\|\psi\|_{H^1(\Ga)}+\|f\|_{H^{-1}(\Om;w)}\right). 
	\end{equation}
	Furthermore, if $f\in L^2(\Om)$, then  $v\in \mcD(A)$ and $v\in H^2(\mcO(\Ga;\be))$. 
\end{lemma}

\begin{equation}\label{09.02.1}
	\begin{cases}
		\pt_{tt}y+A y=f, &\mbox{in }Q,\\
		y(0)=y_0, \pt_ty(0)=y_1, &\mbox{in }\Om,\\
		y=u, &\mbox{on }\Sigma, 
	\end{cases}
\end{equation}

\begin{lemma}\label{09.04.L3}
	Under the conditions in Theorem \ref{09.02.T1} with $u=0$ on $\Sigma$. Then the mapping 
	\begin{equation*}
		H_0^1(\Om;w)\ts L^2(\Om)\ts L^1(0,T; L^2(\Om))\ra L^2(\Sigma), \ (y_0,y_1,f)\mapsto \f{\pt y}{\pt\nu_\mcA}
	\end{equation*}
	is a continuous linear operator, where $y$ is the solution of \eqref{09.02.1} with $u=0$ on $\Sigma$. Furthermore, there exists a constant $C>0$ such that 
	\begin{equation}\label{09.04.6}
		\|\pt_{\nu_\mcA} y\|_{L^2(\Sigma)}\leq C \left(\|y_0\|_{H_0^1(\Om;w)}+\|y_1\|_{L^2(\Om)}+\|f\|_{L^1(0,T; L^2(\Om))}\right).
	\end{equation}
\end{lemma}

\begin{definition}\label{09.04.D1}
	Under Assumption \ref{06.27.A1}, assume that 
	\begin{equation*}
		f=0 \mbox{ in }Q, \quad y_0\in L^2(\Om), \quad y_1\in H^{-1}(\Om;w), \mbox{ and } u\in L^2(\Sigma). 
	\end{equation*}
	We call 
	\begin{equation*}
		y\in C([0,T]; L^2(\Om)), \mbox{ with } \pt_ty \in C([0,T]; H^{-1}(\Om;w)), 
	\end{equation*} 
	is a solution of \eqref{09.02.1}  if,  for each $t\in [0,T]$, and for all  $z_0\in H_0^1(\Om),z_1\in L^2(\Om)$, we have 
	\begin{equation}\label{09.12.1}
		M(z_0,z_1)=(y(t), \pt_t z(t))_{L^2(\Om)}-\lg \pt_ty(t), z(t)\rg_{H^{-1}(\Om;w), H_0^1(\Om;w)}, 
	\end{equation}
	where 
	\begin{equation*}
		M(z_0,z_1)=\int_0^t\int_\Ga  \f{\pt z}{\pt \nu_\mcA}(x,s) u(x,s)\df S\df s+(y_0, z_1)_{L^2(\Om)}-\lg y_1, z_0\rg_{H^{-1}(\Om;w), H_0^1(\Om;w)}, 
	\end{equation*}
	and $z$ is a solution of 
	\begin{equation*}
		\begin{cases}
			\pt_{tt}z+\mcA z=0, &\mbox{in }Q,\\
			z(0)=z_0, \pt_tz(0)=z_1, &\mbox{in }\Om,\\
			z=0, &\mbox{on }\Sigma. 
		\end{cases}
	\end{equation*} 
\end{definition}

\begin{lemma}\label{09.04.L1}
	Under the conditions in Definition \ref{09.04.D1}. Then there exists a unique solution 
	\begin{equation*}
		y\in C([0,T]; L^2(\Om))\cap C^1([0,T]; H^{-1}(\Om;w))
	\end{equation*}
	of \eqref{09.02.1}, and there exists a constant $C>0$, depends only on $w$, such that for $t\in (0,T)$ we have 
	\begin{equation*}
		\|y(t)\|_{L^2(\Om)}+\|\pt_t y(t)\|_{H^{-1}(\Om;w)}\leq C\left(\|y_0\|_{L^2(\Om)}+\|y_1\|_{H^{-1}(\Om;w)}+\|u\|_{L^2(\Sigma)}\right). 
	\end{equation*} 
\end{lemma}

	\begin{equation}\label{08.15.6}
		y=C(t)y_0+S(t)y_1+(Lu)(t), 
	\end{equation}

	\begin{equation}\label{09.03.4}
		\begin{split} 
			\pt_ty 
			&=\mcA S(t) y_0+C(t) y_1+\mcA \int_0^t C(t-\tau)Du(\tau)\df \tau,\\
			(L_tu)(t)
			&=\left(\f{\df L}{\df t}u\right)(t)=\mcA\int_0^t C(t-\tau)Du(\tau)\df\tau. 
		\end{split} 
	\end{equation}

\begin{corollary}\label{09.03.C1}
	Under Assumption \ref{06.27.A1}. Let $u\in H^{1}(\Sigma)$. Then  
	\begin{equation*}
		Lu\in C([0,T]; H^1(\Om;w)),\quad L_tu\in C([0,T]; L^2(\Om)), 
	\end{equation*}
	where $Lu, L_tu$ are defined in \eqref{08.15.6} and \eqref{09.03.4} respectively. 
\end{corollary}

\begin{theorem}\label{09.02.T1}
	Under Assumption \ref{06.27.A1}. Suppose that 
	\begin{equation*}
		f\in L^1(0,T; L^2(\Om)),\quad y_0\in H_0^1(\Om;w),\quad y_1\in L^2(\Om), \mbox{ and } u\in H^1(\Sigma). 
	\end{equation*} 
	Let $y$ be the solution of \eqref{09.02.1}. 	
	Then the solution
	\begin{equation*}
		y\in C([0,T]; H^1(\Om;w))\cap C^1([0,T]; L^2(\Om))
	\end{equation*}
	of \eqref{09.02.1} satisfies: there exists a constant $C>0$, depends only on $w, T$ and $\Om$,  such that 
	\begin{equation}\label{09.04.1}
		\begin{split} 
		&\|y(t)\|_{C([0,T];H^1(\Om;w))}+\|\pt_t y(t)\|_{C([0,T]; L^2(\Om))}\\
		&\leq C\left(\|u\|_{H^1(\Sigma)}+\|y_0\|_{H_0^1(\Om;w)}+\|y_1\|_{L^2(\Om)}+\|f\|_{L^1(0,T; L^2(\Om))}\right).
		\end{split} 
	\end{equation}
	Furthermore, we have 
	\begin{equation*}
		\f{\pt y}{\pt \nu_\mcA}\in L^2(\Sigma), 
	\end{equation*}
	and there exists a constant $C=C(T,\Om)>0$, such that 
	\begin{equation}\label{09.04.2}
		\|\pt_{\nu_\mcA} y\|_{L^2(\Sigma)}\leq C\left(\|u\|_{H^1(\Sigma)}+\|y_0\|_{H^1(\Om;w)}+\|y_1\|_{L^2(\Om)}+\|f\|_{L^1(0,T; L^2(\Om))}\right).
	\end{equation}
\end{theorem}

\begin{proof}[Proof of Theorem \ref{09.02.T1}]
	It is easy to see that the solution of the equation \eqref{09.02.1} can be written as 
	\begin{equation*}
		y=h+z, 
	\end{equation*}
	where $h$ is the solution of the following system
	\begin{equation}\label{09.12.5}
		\begin{cases}
			\pt_{tt}h+\mcA h=f, &\mbox{in }Q, \\
			h(0)=\pt_th(0)=0, &\mbox{in }\Om,\\
			h=0, &\mbox{on }\Sigma, 
		\end{cases}
	\end{equation}
	and $z$ is the solution of the following system
	\begin{equation}\label{09.12.6}
		\begin{cases}
			\pt_{tt}z+A z=0, &\mbox{in }Q, \\
			z(0)=y_0, \pt_tz(0)=y_1, &\mbox{in }\Om,\\
			z=u, &\mbox{on }\Sigma. 
		\end{cases}
	\end{equation}
	
	On one hand, for the system \eqref{09.12.5}, we have Theorem \ref{thm:2.12} and Lemma \ref{09.04.L3}. i.e., 
	\begin{equation*}
		\|h(t)\|_{C([0,T];H^1(\Om;w))}+\|\pt_t h(t)\|_{C([0,T]; L^2(\Om))}\leq C\|f\|_{L^1(0,T; L^2(\Om))}, 
	\end{equation*}
	and 
	\begin{equation*} 
		\|\pt_{\nu_\mcA} h\|_{L^2(\Sigma)}\leq C \|f\|_{L^1(0,T; L^2(\Om))}.
	\end{equation*}
	
	On the other hand, for system \eqref{09.12.6}, we set $\wh z=\pt_tz$, and 
	\begin{equation*}
		\begin{cases}
			\pt_{tt}\wh z+A \wh z=0, &\mbox{in }Q, \\
			\wh z(0)=y_1, \pt_t\wh z(0)=\mcA y_0, &\mbox{in } \Om,\\
			\wh z=\pt_t u, &\mbox{on }\Sigma, 
		\end{cases}
	\end{equation*}
	since $\pt_tu\in L^2(\Sigma), y_1\in L^2(\Om)$, and $\mcA y_0\in H^{-1}(\Om;w)$ by Corollary \ref{09.03.C1} in the following, by Lemma \ref{09.04.L1}, we have 
	\begin{equation*}
		\wh z\in C([0,T]; L^2(\Om))\cap C^1([0,T]; H^{-1}(\Om;w)), 
	\end{equation*}
	and there exists a constant $C>0$ such that for all $t\in (0,T)$, 
	\begin{equation*}
		\|\wh z\|_{L^2(\Om)}+\|\pt_t\wh z(t)\|_{H^{-1}(\Om;w)}\leq C\left(\|y_0\|_{H_0^1(\Om;w)}+\|y_1\|_{L^2(\Om)}+\|u\|_{H^1(\Sigma)}\right)
	\end{equation*}
	by Corollary \ref{09.03.C1} in the following. 
	These imply that 
	\begin{equation*}
		z\in C^1([0,T]; L^2(\Om))\cap C^2([0,T]; H^{-1}(\Om;w)), 
	\end{equation*}
	and from $\pt_{tt}z+A z=0$ we have 
	\begin{equation*}
		Az\in C([0,T]; H^{-1}(\Om;w)). 
	\end{equation*}
	Note that $u(\cdot, t)\in H^1(\Ga)$, from Lemma \ref{09.11.L1}, we have 
	\begin{equation*}
		z\in C([0,T]; H^1(\Om;w))\cap C^1([0,T]; L^2(\Om)). 
	\end{equation*}
	Moreover, there exists $C>0$ such that for all $t\in (0,T)$ we have 
	\begin{equation*}
		\|\pt_tz\|_{L^2(\Om)}+\|Az\|_{H^{-1}(\Om;w)}\leq C\left(\|y_0\|_{H_0^1(\Om;w)}+\|y_1\|_{L^2(\Om)}+\|u\|_{H^1(\Sigma)}\right). 
	\end{equation*}
	Using Lemma \ref{09.11.L1} again, we obtain 
	\begin{equation*}
		\|z(t)\|_{H^1(\Om;w)}+\|\pt_t z(t)\|_{L^2(\Om)}\leq C\left(\|y_0\|_{H_0^1(\Om;w)}+\|y_1\|_{L^2(\Om)}+\|u\|_{H^1(\Sigma)}\right). 
	\end{equation*}
	
	Overall, we have proved Theorem \ref{09.02.T1}. 
\end{proof}

\end{document}
